\begin{abstract}
随着大语言模型（LLM）的快速发展，基于 LLM 的智能体系统在桌面场景中展现出巨大应用潜力。然而，现有桌面智能体方案缺乏结构化的分层架构，意图识别策略单一，视觉感知能力薄弱。本文提出一种面向桌面场景的多模态分层智能体架构，将意图路由、角色扮演和工作流执行解耦为三层独立模块。在路由层，设计了基于关键词组合门控的快速意图分类方法，要求至少两个语义类别的关键词共现才判定为代码任务意图，配合 LLM 结构化参数抽取实现精准操作捕获；在视觉层，基于微调的 YOLOv8 模型构建了桌面 UI 元素检测与活动画像映射管线；在工作引擎层，基于 LangGraph 实现了包含 5 个子图的图工作流编排。在包含 401 条测试用例的实验中，关键词组合门控方法取得了 0.685 的加权 F1 值和 70.1\% 的准确率，路由层延迟仅为 0.020ms。与单类别关键词匹配基线（F1=0.788）相比，组合门控在未知意图类别的 F1 值上提升了 0.374（从 0.039 到 0.413），体现了安全性优先的设计权衡。消融实验验证了组合门控对误分类控制的核心贡献。

\vspace{1em}
\noindent\textbf{关键词：}桌面智能体；意图识别；YOLOv8；LangGraph；角色扮演
\end{abstract}

\begin{abstract}
With the rapid development of large language models (LLMs), LLM-based agent systems have shown great potential in desktop scenarios. However, existing desktop agent solutions lack structured layered architectures, rely on single-strategy intent recognition, and have weak visual perception capabilities. This paper proposes a multi-modal layered desktop agent architecture that decouples intent routing, roleplay, and workflow execution into three independent modules. At the routing layer, we design a keyword-combination gating mechanism requiring co-occurrence of at least two semantic categories for intent classification, paired with LLM-based structured parameter extraction. At the visual layer, we build a desktop UI element detection and activity profiling pipeline based on fine-tuned YOLOv8. At the workflow engine layer, we implement a five-subgraph graph workflow orchestration based on LangGraph. Experiments on a 401-item test set show that our keyword combination gating achieves a weighted F1 of 0.685 and accuracy of 70.1\%, with routing latency of only 0.020ms. Compared to single-category keyword matching (F1=0.788), the combination gate improves the unknown-class F1 by 0.374 (from 0.039 to 0.413), demonstrating a safety-oriented design tradeoff. Ablation studies confirm the contribution of the combination gate to false-positive control.

\vspace{1em}
\noindent\textbf{Keywords:} Desktop Agent; Intent Recognition; YOLOv8; LangGraph; Roleplay
\end{abstract}
